\documentclass[11pt,a4paper]{article}
\usepackage{latex/research-report}
\usepackage{newunicodechar}
\newunicodechar{→}{\ensuremath{\rightarrow}}
\newunicodechar{×}{\ensuremath{\times}}
\begin{document}
\hypersetup{pdftitle={Validation of an Edge-Fog-Cloud Scheduling Testbed},pdfauthor={Gleb Mozhaiskii}}
\section{Validation of an Edge-Fog-Cloud Scheduling
Testbed}\label{validation-of-an-edge-fog-cloud-scheduling-testbed}

\textbf{Gleb Mozhaiskii · Research note · October 2026}

\subsection{Abstract}\label{abstract}

I evaluated whether an application-level scheduling testbed could be
rebuilt and execute the same workload correctly across three fresh Linux
VM sessions. Each session contained a Gateway orchestrator and three
CPU-limited workers connected through MADT, with global, sender and
worker queues. The runtime, 12 scheduling scenarios, validation criteria
and stopping rules were fixed before collection. All three sessions
completed, producing 2,844 reference-checked outputs, including 432
scored tasks. Of these, 431 met the one-second deadline measured from
planned arrival; one correct result was late. The configured network
conditions, evidence export and cleanup also passed their checks. The
results support repeatable execution of this workload on the tested
host. Placement, latency and decision overhead are reported as
exploratory observations.

\subsection{1. Research Objective}\label{research-objective}

The predefined question was whether the existing orchestrator could
reconstruct its runtime and correctly execute and account for the
original pilot across \textbf{three fresh VM sessions on the same host}.
A separate question concerned reproduction of the configured network
conditions within the original numerical tolerances.

The application screen required three complete sessions. Completion
required all expected task outcomes and reference outputs, consistent
assignments and queue transitions, one execution slot per worker, finite
resource accounting, valid clocks and provenance, intact evidence
export, and verified cleanup. Network qualification was assessed
independently. Strict performance eligibility additionally required both
network qualifications to pass and the preserved evidence and cleanup to
be verified. A correct result that missed its deadline remained a late
result; lateness alone did not fail application correctness.

RR, Compute-ECT and Network-ECT exercised the instrument under different
placement decisions. Their performance ordering was not a validation
criterion. The policy comparisons in Section 6 are exploratory
observations from the completed cohort.

Wang et al. {[}1{]} provide methodological context: compare controllers
under common conditions and report decision overhead alongside
service-level outcomes. Their Continuous Multi-Mode Scheduling benchmark
selects execution modes as well as destinations for video inference.
This testbed instead assigns fixed CPU tasks to workers and checks their
results.

\Needspace{100pt}

\subsection{2. Testbed Architecture}\label{testbed-architecture}

Each session used a fresh ARM64 Lima Linux VM with 4 vCPU, 4 GiB RAM and
a 20 GiB sparse disk. The three sessions ran sequentially on the same
Apple Silicon Mac. MADT {[}2{]} deployed the four-container topology
below. G/E/F/C CPU quotas were 0.50/0.25/0.50/1.00 CPU, with a 100 ms
period and a 256 MiB memory limit per container.

\begin{figure}[H]
\centering
\includegraphics[width=1\linewidth,height=\textheight,keepaspectratio,alt={System topology}]{build/figures/topology.png}
\caption{System topology}
\end{figure}

Gateway G owns the global FIFO, placement controller and one sender FIFO
per worker. Each sender serializes uploads until input-admission
acknowledgement. Workers E, F and C each have a local FIFO and one
non-preemptive execution slot; execution can overlap the next input
transfer. G reaches C through F, while E is directly reachable on the
access network. F therefore shares its quota between routing and
application work.

Each worker permits at most 32 unresolved assignments and 64 MiB of
application reservations. A task reserves its input size plus 1 MiB;
capacity is released after a validated result. These bounds are distinct
from the container memory limit. Worker state is reported every 100 ms,
and placement uses an immutable snapshot of controller-known events and
state.

The workload is deterministic PBKDF2-HMAC-SHA256 computation over a
digest of the transferred input. Two input sizes (64 KiB and 1 MiB) and
two work levels (200,000 and 800,000 iterations) define four task
classes. G transfers the complete input and independently checks the
returned result vector. Task, assignment and output identities are
checked within each execution group.

\subsection{3. Scheduling Policies}\label{scheduling-policies}

All policies require an online worker with a free assignment credit and
enough reservation capacity. A busy CPU remains eligible because work
can wait in its queues.

\textbf{Round Robin (RR)} cycles through E, F and C, skipping ineligible
workers. It provides a baseline without completion-time prediction.

\textbf{Compute-ECT} minimizes predicted completion time using the
current task\textquotesingle s remaining execution estimate, assigned
waiting work and the new task\textquotesingle s calibrated service time.
Transport estimates are deliberately set to zero, providing a
compute-only control with the same queue accounting.

\textbf{Network-ECT} adds calibrated input delivery, admission
acknowledgement and result-return costs. It tracks sender serialization
separately from worker execution, allowing the next input transfer to
overlap the current computation. Result return extends the response
estimate without occupying the worker\textquotesingle s execution slot.

Both ECT policies use only controller-known state and calibration
profiles frozen before comparison, with fixed E/F/C tie-breaking. The
exact equations, profile construction and observation limits are
documented in \texttt{docs/scheduling-model.md} in the accompanying
repository.

\Needspace{120pt}

\subsection{4. Experimental Protocol}\label{experimental-protocol}

The prospective cohort comprised three sequential fresh VM sessions, Run
1-3, attempted once each without replacement or tuning. Each
reconstructed the pinned runtime and used the same task trace, policy
definitions, quotas and hash-derived scenario order. The baseline
included Lima 2.2.1, Docker 28.5.2, containerd 1.7.29 and the pinned
MADT revision and network-control repair. Full source identities,
dependency pins and execution commands accompany the public artifact.

Each session crossed three policies, two arrival patterns and two
network conditions: \textbf{12 scenarios, each with 12 scored offers
over six seconds}. The steady trace releases one task every 0.5 s. The
burst trace combines a one-task/s stream with three additional tasks
near 2 s and three near 4 s. Task classes and release times are
identical across policies and sessions.

The scheduling conditions were unshaped networking and a nominal 20 ms
one-way delay on Cloud ingress from Fog. The separate 10 Mbps treatment
qualified network control; it was not a scored scheduling condition.
Each VM performed 72 calibration tasks, 12 queue-check tasks, 720 common
warmup tasks and 144 scored tasks: 948 executions in total.

\textbf{Application criteria.} All 15 execution groups, including the 12
scored scenarios, had to complete with 948 independently checked worker
outputs and all 144 scored outcomes retained. The audit checked task
conservation, assignment/output identities, FIFO and single-slot
constraints, finite credits and reservations, queue/execution overlap,
causal timestamps, resource limits and source/runtime provenance.
Missing, invalid or incompletely audited application evidence required a
stop.

\textbf{Network criteria.} Both health and application-quota
qualification measured baseline, impaired and restored states. For both
probe directions, the median RTT increment had to be 15-25 ms and the
restored median within 2 ms of baseline. The unaffected access-path
medians could span at most 2 ms. The rate treatment required 7-11 Mbps,
with baseline and restoration at least 20 Mbps and restoration at least
75\% of baseline. Ordinary measurements required 20/20 ping delivery and
valid TCP delivery. Route, interface, qdisc and numerical actuation
readbacks, partition/recovery and unaffected-path checks were also
mandatory. A valid numerical deviation was recorded separately from
application correctness and excluded the session from strict performance
comparison; it did not become a relaxed network pass.

\textbf{Measurement and preservation.} A scored task meets its deadline
when \(z_i-a_i^*\leq1\) s, from planned release \(a_i^*\) to verified
response \(z_i\). The record separates release lag, global waiting,
observation, decision, commitment, sender waiting, input delivery,
worker waiting, execution and result return. Decision wall time measures
the timed policy call; thread CPU time is recorded separately.
Containers share a checked Linux monotonic-clock domain. Export
integrity and actual session cleanup were verified before publication
eligibility was finalized.

\Needspace{200pt}

\subsection{5. Validation Results}\label{validation-results}

All three sessions passed the application screen and both network
qualifications. Evidence export, preservation and cleanup passed their
recorded audits. All 432 scored offers completed correctly: none was
rejected, failed, unresolved or unreleased. The three sessions produced
2,844 reference-checked worker outputs including preparation.

\begin{table}[H]
\centering
\label{table:sessions:start}
\begin{tabular}{llrr}
\toprule
\textbf{Session} & \textbf{Application} & \textbf{Correct outputs} & \textbf{Scored on time} \\
\midrule
Run 1 & Complete & 948 & 144/144 \\
Run 2 & Complete & 948 & 144/144 \\
Run 3 & Complete & 948 & 143/144 \\
\bottomrule
\end{tabular}
\label{table:sessions:end}
\end{table}


The Gateway recorded 948 assignments per session, including 144 scored
decisions. Worker totals below include calibration, queue checks and
warmup; parentheses contain scored outputs. Every worker produced
verified results in each session.

\begin{table}[H]
\centering
\label{table:nodes:start}
\begin{tabular}{lrrrr}
\toprule
\textbf{Session} & \textbf{G all / scored} & \textbf{Edge} & \textbf{Fog} & \textbf{Cloud} \\
\midrule
Run 1 & 948 / 144 & 284 (16) & 294 (26) & 370 (102) \\
Run 2 & 948 / 144 & 284 (16) & 294 (26) & 370 (102) \\
Run 3 & 948 / 144 & 284 (16) & 298 (30) & 366 (98) \\
\bottomrule
\end{tabular}
\label{table:nodes:end}
\end{table}


The application-quota network qualification produced the following RTT
increments and limited receive rates. Forward and reverse denote probe
initiation directions; these are RTT increments, not direct
one-way-delay estimates. All values lie within the original 15-25 ms and
7-11 Mbps bands; restoration, delivery and control-path checks also
passed.

\begin{table}[H]
\centering
\label{table:network:start}
\begin{tabular}{lrrr}
\toprule
\textbf{Session} & \textbf{Forward delta (ms)} & \textbf{Reverse delta (ms)} & \textbf{Limited rate (Mbps)} \\
\midrule
Run 1 & 22.139 & 20.779 & 9.072 \\
Run 2 & 21.223 & 21.570 & 9.320 \\
Run 3 & 22.045 & 21.967 & 8.357 \\
\bottomrule
\end{tabular}
\label{table:network:end}
\end{table}


Of the 432 scored tasks, \textbf{431 met the one-second planned-arrival
deadline}. The remaining task was an Edge assignment in Run 3,
RR/burst/unshaped. Its correct response arrived after 1002.317940 ms,
exceeding the deadline by 2.317940 ms. Release lag was 4.019139 ms;
measuring from actual release would yield 998.298801 ms and change its
classification. The predefined planned-arrival result is retained.

\begin{figure}[H]
\centering
\includegraphics[width=1\linewidth,height=\textheight,keepaspectratio,alt={Run 3 task timeline for Round Robin with burst arrivals; the late task is retained}]{build/figures/task-timeline.pdf}
\caption{Run 3 task timeline for Round Robin with burst arrivals; the
late task is retained}
\end{figure}

\Needspace{500pt}

\subsection{6. Exploratory Observations}\label{exploratory-observations}

The \textbf{VM session is the replication unit: N=3 on one host}. Each
run cell below gives median response time in milliseconds and, beneath
it, assignments to Edge / Fog / Cloud, out of 12 tasks.
\textquotesingle On time\textquotesingle{} is a descriptive total over
all 36 offers in that row. The rows and tasks are nested observations,
not additional independent replications. All 12 scenario combinations
are shown.

\begin{table}[H]
\centering
\label{table:scenario-summary:start}
\begin{tabular}{lllrrrr}
\toprule
\textbf{Arrival} & \textbf{Network} & \textbf{Scheduler} & \textbf{Run 1} & \textbf{Run 2} & \textbf{Run 3} & \textbf{On time} \\
\midrule
Steady & Normal & Round Robin & \shortstack[r]{181.7\\4 / 4 / 4} & \shortstack[r]{236.7\\4 / 4 / 4} & \shortstack[r]{227.2\\4 / 4 / 4} & 36/36 \\
Steady & Normal & Compute-ECT & \shortstack[r]{118.0\\0 / 0 / 12} & \shortstack[r]{121.2\\0 / 0 / 12} & \shortstack[r]{119.2\\0 / 0 / 12} & 36/36 \\
Steady & Normal & Network-ECT & \shortstack[r]{113.6\\0 / 0 / 12} & \shortstack[r]{120.1\\0 / 0 / 12} & \shortstack[r]{119.2\\0 / 0 / 12} & 36/36 \\
Steady & +20 ms & Round Robin & \shortstack[r]{266.1\\4 / 4 / 4} & \shortstack[r]{252.4\\4 / 4 / 4} & \shortstack[r]{248.9\\4 / 4 / 4} & 36/36 \\
Steady & +20 ms & Compute-ECT & \shortstack[r]{209.0\\0 / 0 / 12} & \shortstack[r]{203.5\\0 / 0 / 12} & \shortstack[r]{199.1\\0 / 0 / 12} & 36/36 \\
Steady & +20 ms & Network-ECT & \shortstack[r]{196.3\\0 / 3 / 9} & \shortstack[r]{207.0\\0 / 3 / 9} & \shortstack[r]{145.7\\0 / 6 / 6} & 36/36 \\
Burst & Normal & Round Robin & \shortstack[r]{185.5\\4 / 4 / 4} & \shortstack[r]{171.9\\4 / 4 / 4} & \shortstack[r]{189.8\\4 / 4 / 4} & 35/36 \\
Burst & Normal & Compute-ECT & \shortstack[r]{131.2\\0 / 1 / 11} & \shortstack[r]{143.3\\0 / 1 / 11} & \shortstack[r]{136.0\\0 / 1 / 11} & 36/36 \\
Burst & Normal & Network-ECT & \shortstack[r]{137.3\\0 / 1 / 11} & \shortstack[r]{138.3\\0 / 1 / 11} & \shortstack[r]{139.0\\0 / 1 / 11} & 36/36 \\
Burst & +20 ms & Round Robin & \shortstack[r]{210.6\\4 / 4 / 4} & \shortstack[r]{210.4\\4 / 4 / 4} & \shortstack[r]{210.2\\4 / 4 / 4} & 36/36 \\
Burst & +20 ms & Compute-ECT & \shortstack[r]{191.0\\0 / 1 / 11} & \shortstack[r]{186.8\\0 / 1 / 11} & \shortstack[r]{198.7\\0 / 1 / 11} & 36/36 \\
Burst & +20 ms & Network-ECT & \shortstack[r]{186.9\\0 / 4 / 8} & \shortstack[r]{206.5\\0 / 4 / 8} & \shortstack[r]{204.5\\0 / 5 / 7} & 36/36 \\
\bottomrule
\end{tabular}
\label{table:scenario-summary:end}
\end{table}


\begin{figure}[H]
\centering
\includegraphics[width=1\linewidth,height=\textheight,keepaspectratio,alt={Median task response time in each fresh VM run}]{build/figures/median-latency.pdf}
\caption{Median task response time in each fresh VM run}
\end{figure}

Each symbol is one session\textquotesingle s scenario median. The
horizontal segment spans the observed minimum and maximum of the three
medians; it is not a confidence interval.

\Needspace{55pt}

Both ECT response medians were lower than the matched RR median in every
observed session/context pair. Across fixed scenarios, the relative
range \((\max-\min)/\mathrm{median}\) was 0.219\%-31.267\%.

\textbf{Placement.} RR assigned 4/4/4 in every scenario. Neither ECT
policy selected Edge in this cohort. With burst arrivals and added
delay, Compute-ECT assigned one task to Fog per run; Network-ECT
assigned 4, 4 and 5. Under steady arrivals and added delay, Network-ECT
assigned 3, 3 and 6 tasks to Fog. Its corresponding latency medians were
196.273, 207.037 and 145.668 ms, the largest inter-session spread. The
placement change and lower Run 3 median are associated observations;
calibration and execution timing also varied between sessions.

\textbf{Decision cost and offered load.} Decision wall medians across
the 36 scenario/session points ranged from 9.375 to 160.792
microseconds; CPU medians ranged from 6.459 to 146.083 microseconds. In
all 24 matched ECT-minus-RR comparisons, both decision wall and CPU
medians were higher for ECT. Decision cost remains included in measured
task response.

\begin{figure}[H]
\centering
\includegraphics[width=1\linewidth,height=\textheight,keepaspectratio,alt={Scheduling decision wall time in all three sessions; each point is one scenario median}]{build/figures/decision-overhead.pdf}
\caption{Scheduling decision wall time in all three sessions; each point
is one scenario median}
\end{figure}

Every scenario completed its 12 scored tasks within the six-second
window: observed window throughput was 2 tasks/s. On-time goodput was
also 2 tasks/s except for the late-task scenario, where it was 11/6
tasks/s. This offered load does not identify maximum capacity or
overload behavior.

\subsection{7. Threats to Validity}\label{threats-to-validity}

\textbf{Replication and scope.} Three fresh VMs shared one physical host
and operator. They tested runtime reconstruction and repeated execution
with one recipe, a fixed seed and a fixed scenario order. The data do
not establish portability, long-term reliability or a general policy
ranking. CPU quotas and emulated links represent logical tiers.

\textbf{Timing and attribution.} CPU contention, quota replenishment,
virtualization, network emulation and measurement activity all
contribute to wall time. They were not varied independently, so the
deadline miss and inter-session spread cannot be assigned to a single
cause. The range of three medians is descriptive uncertainty, with too
little replication for broad statistical conclusions. Decision wall time
includes elapsed scheduling and instrumentation effects; it is not total
control-plane cost.

\textbf{Workload and models.} Twelve scored tasks per scenario and four
deterministic CPU classes provide a correctness workload. They do not
cover stateful services, application dependency graphs, migration,
recovery of interrupted work or sustained saturation. Service profiles
come from unshaped calibration; transport profiles are
condition-specific and fixed. Model error under changing load or stale
state was not experimentally isolated. Fog also couples forwarding and
execution.

\Needspace{65pt}

\textbf{Reconstruction.} Source, versions and package manifests were
pinned and checked, but the environment was not rebuilt from a hermetic
package mirror. Future dependency availability remains an operational
risk. All three collected sessions are included; preparation and warmup
outputs are distinguished from scored outcomes.

\subsection{8. Conclusion and Research
Directions}\label{conclusion-and-research-directions}

The predefined same-host application screen passed in three fresh VM
sessions, with complete task accounting and verified outputs on every
worker. The configured network treatments passed their separate
qualifications, and preserved evidence and cleanup supported final
comparison eligibility. The scored outcome was 431 on-time results and
one late correct result under the original planned-arrival deadline.

A proposed follow-up would introduce controlled delays in worker-state
updates while holding the workload, policies and configured network
conditions fixed. It would measure completion-time prediction error,
task placement, planned-arrival deadline attainment and the CPU and
traffic costs of state updates. Varying update delay and frequency
separately would help distinguish the effect of stale information from
the cost of collecting it. This would test a research question; the
current cohort does not establish a stale-telemetry defect. Temporary
worker loss and recovery would require a separate failure and task-state
model.

\subsection{References}\label{references}

\begin{enumerate}
\def\labelenumi{\arabic{enumi}.}
\tightlist
\item
  Z. Wang, B. Sedlak, J. L. Herrera, and S. Dustdar.
  \textquotesingle Fair Comparison of Scheduling Algorithms on
  Heterogeneous Edge Clusters: A Continuous Adaptive
  Benchmark.\textquotesingle{} IEEE ICDCSW, pp. 232-237, 2026.
  \href{https://doi.org/10.1109/ICDCSW72724.2026.00054}{DOI:
  10.1109/ICDCSW72724.2026.00054}. Open preprint:
  \href{https://arxiv.org/abs/2606.12343}{arXiv:2606.12343}.
\item
  MADT project. \textquotesingle Distributed Application Modelling
  System.\textquotesingle{}
  \href{https://github.com/dltcspbu/madt}{Source repository}, accessed
  23 October 2026.
\end{enumerate}

\end{document}
